\documentclass[11pt, a4paper]{article}
\usepackage[a4paper, top=2.5cm, bottom=2.5cm, left=2cm, right=2cm]{geometry}





% Add because main language is not English
\usepackage{enumitem}
\setlist[itemize]{label=-}
% ---------------------------------

\usepackage{amsmath}
\usepackage{amssymb}
\usepackage{booktabs}
\usepackage{tabularx}
\usepackage{listings}
\usepackage{xcolor}

\definecolor{codegreen}{rgb}{0,0.6,0}
\definecolor{codegray}{rgb}{0.5,0.5,0.5}
\definecolor{codepurple}{rgb}{0.58,0,0.82}
\definecolor{backcolour}{rgb}{0.95,0.95,0.92}

\oddsidemargin -1.5cm 
\textwidth 19 cm

\pagestyle{empty}

\begin{document}


\begin{center}
	\LARGE \textbf{Presentaciones Orales de acuerdo a la \\ GUÍA DE APRENDIZAJE
		PROGRAMACIÓN. \\ GRADO EN INGENIERÍA AEROESPACIAL.} \\
	\LARGE \textbf{Rúbrica de Evaluación Docente.}
\end{center}
\vspace{0.5cm}

\noindent
\begin{tabularx}{\textwidth}{|p{1cm}|p{6.5cm}|p{10.2cm}|}
	\toprule
	\textbf{Peso} & \textbf{Categoría de Evaluación} & \textbf{Criterios de Excelencia}  \\
	\midrule
	\textbf{30\%} & \textbf{Selección de Métodos Numéricos e Implementación Matemática} \newline \small Evalúa la implementación matemática, el uso de \texttt{numpy} / \texttt{scipy} y la exactitud de los resultados numéricos en los 6 apartados. 
	& \begin{itemize}[leftmargin=*, nosep]
		\item Resuelve correctamente el sistema lineal.
		\item Realiza la regresión adecuadamente.
		\item Interpola correctamente.
		\item Implementa diferencias finitas sin errores en índices o extremo.
		\item Integra correctamente aplicando Simpson/Trapecios.
		\item Encuentra la raíz  con tolerancia adecuada.
	\end{itemize} \\
	\midrule
	\textbf{30\%} & \textbf{Calidad y Estructura del Código} \newline \small Evalúa las buenas prácticas de programación, modularización, paradigma funcional, vectorización frente a bucles for, y nomenclatura.
	& \begin{itemize}[leftmargin=*, nosep]
		\item Código modular, uso adecuado de \texttt{def} para separar lógica y funciones del GUI.
		\item Excelente uso de la vectorización de \texttt{numpy} (sin bucles \texttt{for} innecesarios en arrays).
		\item Nomenclatura descriptiva.
		\item Código limpio, ordenado y fácil de leer.
	\end{itemize}  \\
	\midrule
	\textbf{10\%} & \textbf{Interpretación Física} \newline \small Evalúa la capacidad del alumno para conectar el número con la realidad del problema aeroespacial.
	& \begin{itemize}[leftmargin=*, nosep]
		\item Las salidas detallan claramente las magnitudes con sus unidades físicas (SI).
		\item Contiene comentarios que validan el orden de magnitud de los resultados.
		\item Justifica la elección de puntos iniciales en el cálculo de raíces.
	\end{itemize}  \\
	\midrule
	\textbf{10\%} & \textbf{Visualización de Datos} \newline \small Evalúa el dominio de \texttt{matplotlib} y \texttt{matplotlib.widgets} para crear una GUI interactiva.
	& \begin{itemize}[leftmargin=*, nosep]
		\item Interfaz única con navegación funcional.(\texttt{Buttons}) 
		\item Implementa un \texttt{Slider} que recalcula y actualiza una gráfica dinámicamente.
		\item Integración de los resultados numéricos como texto dentro de la GUI.
		\item Gráficas con calidad científica: títulos, etiquetas con unidades, grid y leyendas.
	\end{itemize}  \\
		\midrule
	\textbf{20\%} & \textbf{Claridad de la exposición} 
	\newline \small Evalúa un discurso articulado, fluido y con un hilo conductor nítido. Se pondera la precisión del vocabulario técnico, la coherencia lógica en la transición entre las ideas expuestas.
	& \begin{itemize}[leftmargin=*, nosep]
		\item Existe un hilo argumental de principio a fin.
		\item El vocabulario utilizado es riguroso, técnico y adecuado al nivel del curso. 
		\item Soporte visual minimalista, coherente y en fase con la voz del ponente (las diapositivas no contienen bloques de texto que compitan con la voz del ponente).
		\item Ausencia de muletillas. El discurso tiene fluidez natural.
	\end{itemize}  \\
	\bottomrule
\end{tabularx}

\end{document}

